\chapter{Discussion}

The studies presented in this dissertation examine different contributions of AI to the Climate Services value chain. Their findings raise three central questions: whether high average performance translates into reliability under specific local conditions; how optimization objectives relate to scientific and operational requirements; and whether AI can provide a reliable interface to external climate knowledge. These questions concern not only the capabilities of individual models, but also the conditions under which their outputs can support Climate Services.

\section{Average Skill and Local Reliability}

Several results reveal a distinction between reproducing broad climatic regularities and representing the departures from those regularities that matter locally. This is particularly evident in the LLM climate-change experiment. The models reproduce an approximately plausible mean warming, but the spatial correlations between the warming fields generated by GPT-5 and Mistral Small 3.1 24B and the ERA5 reference remain below 0.10. Their estimates therefore contain substantially less useful geographical information than their plausible average values might suggest.

The weather-model experiments reveal a related, although distinct, problem. Models initialized from EERIE-AMIP initially move towards an ERA5-like global mean, while retaining regional differences that are not apparent in the global average. This adjustment should not be equated with the compression of spatial variability observed in the LLM estimates. Nevertheless, it raises a similar question about the ability to preserve conditions that differ from the regularities represented during training. Closer agreement with a training reference does not, by itself, demonstrate preservation of the climate state supplied at initialization.

These behaviours motivate a broader question: \emph{are models optimized for high average performance also reliable under the highly specific conditions that matter to Climate Services?} Training objectives determine which errors receive greater weight, while the training data determine which conditions are represented. For deterministic predictions, squared-error minimization favours the conditional mean, potentially smoothing uncertain features. Other objectives and architectures behave differently, and the experiments presented here do not establish a common causal mechanism across models. They do, however, show why aggregate performance is insufficient evidence of local reliability.

This distinction is consequential when information is used near an operational threshold. An error that contributes little to a global score may change whether a condition is classified as hazardous at a particular location. Evaluation should therefore examine the geographical distribution, seasonality, and conditional structure of errors, rather than treating local applicability as a consequence of good average skill.

\section{From Metric Optimization to Scientific Relevance}

A second issue concerns the relationship between what an algorithm optimizes and what a Climate Service needs. A performance metric formalizes a criterion of success, but its improvement has scientific value only insofar as that criterion represents the intended task. Contextual relevance is not automatically contained in an optimization objective; it must be introduced through problem formulation, data selection, methodological choices, and evaluation.

The seasonal framework illustrates this distinction. The algorithm does not independently determine what a meteorological season should mean. The input variables, normalization, temporal processing, and number of classes establish an operational representation of the concept. Its scientific interpretation is then assessed against climatological knowledge and recognizable seasonal behaviour. Machine Learning makes this representation computable and transferable, but does not remove the assumptions through which it was constructed.

Consequently, the projected lengthening of the monsoon describes increased persistence of the multivariate conditions associated with the learned seasonal class. It does not automatically imply greater rainfall totals, improved water availability, or a longer agricultural growing season. These interpretations require additional evidence. The relevant question is therefore not simply whether classification performance improves, but whether the resulting information represents the phenomenon required by the application.

Human expertise is fundamental at this stage, before any final output is reviewed. It determines which distinctions are scientifically meaningful and which conclusions an analysis can support. Automation can reduce the effort required to implement these choices, but cannot establish their contextual appropriateness through optimization alone. Expert involvement should therefore be understood as part of the design of AI-assisted Climate Services, rather than only as a safeguard against occasional model failures.

\section{AI as an Interface to Scientific Knowledge}

The limitations of parametric climate recall suggest a different role for language models: mediating access to scientific evidence rather than substituting for its sources. In this configuration, datasets, publications, and analytical software provide the evidential basis, while the model supports question interpretation, information selection, analysis configuration, and communication.

The evaluation problem consequently changes. It concerns how faithfully the system retrieves, transforms, and communicates the information available to it. Relevant questions include whether the selected evidence addresses the user request, whether computational parameters preserve its intended meaning, and whether the final response distinguishes calculated results from interpretation. An answer can contain correct numerical values while misrepresenting their implications.

The augmentation experiments show that these capabilities depend on system design. Contextual augmentation improves specificity and practical orientation without consistently improving perceived credibility or uncertainty communication. Systems incorporating indicator computation show more consistent gains across the evaluated dimensions. This supports the contribution of \emph{procedural grounding}, through which quantitative statements become associated with explicit data, indicator definitions, parameters, and calculations.

Procedural grounding nevertheless establishes traceability, not automatic appropriateness. A reproducible calculation may use an unsuitable threshold, and a correctly computed indicator may not justify a particular adaptation recommendation. The scientific quality of the interface therefore depends both on preserving the evidence it receives and on respecting the limits of the conclusions that evidence supports.

\section{From Output Quality to Decision Support}

The distinction between a better response and a better decision remains unresolved by the experiments presented here. Higher relevance, credibility, or actionability scores indicate improvements according to the adopted evaluation criteria; they do not demonstrate that users make more effective adaptation decisions. Similarly, stronger agreement between human and LLM evaluators after aggregation supports comparative assessment within the tested systems, but does not certify individual recommendations.

Operational evaluation should therefore distinguish analytical correctness, communication quality, and practical usefulness. These dimensions require complementary evidence. Component-level checks can identify errors in data selection or computation, while collaborative evaluations with providers and users are needed to assess whether the information addresses an actual decision problem. Such evaluations should also examine whether explicit uncertainty and appropriate qualification help users interpret results, rather than rewarding confidence independently of its evidential basis.

\section{Scope and Evolution of the Proposed Frameworks}

The findings concern selected models, configurations, variables, and applications. They should not be interpreted as immutable limitations of AI or as evidence that all architectures exhibit the same behaviour. In particular, the weather-model analysis remains an initial assessment, while the knowledge-oriented experiments do not constitute validation of a complete operational Climate Service.

The frameworks developed here seek to separate scientific requirements from specific implementations, making their approach as model-agnostic as possible. This independence is necessarily partial. New training objectives, architectures, and analytical capabilities may alter existing failure modes and introduce others that the present benchmarks do not capture. Evaluation methods must therefore remain open to revision: a framework that is appropriate for current systems may be insufficient for future developments.

The more durable contribution lies in identifying what must be examined: performance under relevant conditions, the scientific meaning of optimized objectives, traceability of evidence, and the validity of its interpretation. These requirements provide a basis for integrating evolving AI capabilities into Climate Services while retaining scrutiny of the assumptions and judgments on which their usefulness depends.